\documentclass[11pt,a4paper]{article}
\usepackage[utf8]{inputenc}
\usepackage[T1]{fontenc}
\usepackage[brazil]{babel}
\usepackage[a4paper,margin=2.3cm]{geometry}
\usepackage{enumitem}
\usepackage{hyperref}
\usepackage{xcolor}
\setlength{\parindent}{0pt}
\setlength{\parskip}{5pt}
\sloppy
\setlength{\emergencystretch}{3em}
\hypersetup{colorlinks=true,urlcolor=blue,linkcolor=black}
\newcommand{\cabecalho}[2]{\begin{center}{\small \textbf{Sprint 2 -- #1}}\\[1mm]\textit{Inventário de Ativos e Vulnerabilidades}\\Matheus Moreli Durão Silva -- 12621CBS217\\Cibersegurança -- Bacharelado -- Noturno\end{center}\hrule\vspace{4mm}}
\newcommand{\bloco}[1]{\par\vspace{2mm}\textbf{#1}\par}
\newcommand{\codigo}[1]{\texttt{#1}}

\begin{document}\cabecalho{S2\_10 -- Atualização de Ativos}{}
\bloco{Regra principal} O ID é a identidade do ativo. Por isso \codigo{atualizar\_ativo} recebe o ID apenas para localizar a chave e não possui parâmetro para trocá-lo. Os campos permitidos são nome, responsável, localização, tipo e criticidade.
\bloco{Como funciona} Valores vazios significam ``manter o anterior'' na tela. Valores novos passam pelas mesmas validações do cadastro. Se houver alteração, \codigo{registrar\_historico} acrescenta data/hora e os campos alterados, criando evidência verificável.
\bloco{Exemplo} Alterar responsável de ``Equipe TI'' para ``Segurança'' muda somente esse atributo. A consulta continua exibindo ID 10. O histórico registra que o responsável mudou.
\bloco{Como testar} Atualizar o ativo 10, conferir o novo valor, conferir ID 10 e verificar que o histórico possui ao menos o evento de criação e o de atualização. Atualizar ID ausente retorna mensagem controlada.
\bloco{Checklist atendido} Identificador preservado, campos limitados, valores validados, inexistência tratada e histórico como evidência.
\end{document}
